% ═══════════════════════════════════════════════════════════════════════
%  APPENDIX A — EVALUATION DATASET
% ═══════════════════════════════════════════════════════════════════════
\section*{Appendix A: Evaluation Dataset}
\addcontentsline{toc}{section}{Appendix A: Evaluation Dataset}

\hs The complete evaluation dataset consists of 30 questions written in natural Khmer,
each annotated with the article(s) of law that correctly answer it and the legal
category it belongs to. The dataset was constructed manually by reading the source laws
and formulating realistic questions a member of the public might ask.

\begin{longtable}{p{0.5cm} p{7.5cm} p{2.2cm} p{3cm}}
\toprule
\textbf{\#} & \textbf{Question (Khmer)} & \textbf{Expected Article(s)} & \textbf{Category} \\
\midrule
\endhead
1  & {\khmerfont តើកិច្ចសន្យាការងារមានអ្វីដែរ?} & 65 & Employment Contract \\ \midrule
2  & {\khmerfont តើកិច្ចសន្យាការងារមានថិរវេលាកំណត់អាចមានរយៈពេលអតិបរមាប៉ុន្មាន?} & 67 & Employment Contract \\ \midrule
3  & {\khmerfont តើរយៈពេលសាកល្បងការងារអតិបរមាមានប៉ុន្មាន?} & 68 & Employment Contract \\ \midrule
4  & {\khmerfont តើរបៀបឱ្យដំណឹងមុនពេលបញ្ឈប់ការងារមានថិរវេលាប៉ុន្មាន?} & 75 & Termination \\ \midrule
5  & {\khmerfont តើការបញ្ឈប់ការងារដោយគ្មានការជូនដំណឹងមុននឹងប៉ះពាល់ដល់និយោជកដោយរបៀបណា?} & 77 & Termination \\ \midrule
6  & {\khmerfont តើម៉ោងធ្វើការធម្មតាប្រចាំថ្ងៃ និងប្រចាំសប្តាហ៍មានប៉ុន្មាន?} & 137 & Working Hours \\ \midrule
7  & {\khmerfont តើម៉ោងបន្ថែមគួរទទួលបានប្រាក់ឈ្នួលប៉ុន្មានភាគរយ?} & 139 & Working Hours \\ \midrule
8  & {\khmerfont តើការងារពេលយប់ ត្រូវបើកប្រាក់ឈ្នួលក្នុងអត្រាប៉ុន្មាន?} & 144 & Working Hours \\ \midrule
9  & {\khmerfont តើការឈប់សម្រាកប្រចាំសប្តាហ៍ត្រូវមានរយៈពេលអប្បបរមាប៉ុន្មាន?} & 147, 146 & Leave \& Holidays \\ \midrule
10 & {\khmerfont តើការឈប់សម្រាកប្រចាំឆ្នាំដោយមានប្រាក់ឈ្នួលមានប៉ុន្មានថ្ងៃ?} & 166 & Leave \& Holidays \\ \midrule
11 & {\khmerfont តើថ្ងៃបុណ្យជាតិ ត្រូវរាប់ជាថ្ងៃឈប់សម្រាកប្រចាំឆ្នាំដែរទេ?} & 166, 161 & Leave \& Holidays \\ \midrule
12 & {\khmerfont តើប្រាក់ឈ្នួលត្រូវបង់ជូនកម្មករនិយោជិតពេលណា មុនការឈប់សម្រាក?} & 168 & Leave \& Holidays \\ \midrule
13 & {\khmerfont តើសហគ្រាសណាខ្លះដែលត្រូវតែមានបទបញ្ជាផ្ទៃក្នុង?} & 22 & Workplace Regulations \\ \midrule
14 & {\khmerfont តើនិយោជកអាចដាក់ទោសវិន័យដល់កម្មករក្នុងរយៈពេលដប់ប្រាំថ្ងៃ ឬច្រើនជាងនេះ?} & 26 & Workplace Regulations \\ \midrule
15 & {\khmerfont តើច្បាប់ការងារហាមឃាត់ការរើសអើងដោយហេតុអ្វីខ្លះ?} & 12 & Discrimination \\ \midrule
16 & {\khmerfont តើប្រាក់បំណាច់បញ្ចប់កិច្ចសន្យាការងារ ត្រូវគណនាដោយរបៀបណា?} & 73 & Wages \& Compensation \\ \midrule
17 & {\khmerfont តើបទបញ្ជាផ្ទៃក្នុងត្រូវកើតឡើងក្នុងរយៈពេលប៉ុន្មានបន្ទាប់ពីបើកសហគ្រាស?} & 24 & Employment Contract \\ \midrule
18 & {\khmerfont តើស្ថាប័នណាទទួលបន្ទុកជំនុំជម្រះវិវាទការងារ?} & 387, 389 & Dispute Resolution \\ \midrule
19 & {\khmerfont តើត្រូវការសមាជិកអប្បបរមាប៉ុន្មាននាក់ ដើម្បីបង្កើតសហជីពមូលដ្ឋាន?} & 10 & Trade Union \\ \midrule
20 & {\khmerfont តើកម្មករនិយោជិតមានសិទ្ធិអ្វីខ្លះទាក់ទងនឹងការចូលរួមសហជីព?} & 5, 7 & Trade Union \\ \midrule
21 & {\khmerfont តើការចុះបញ្ជីសហជីពត្រូវរង់ចាំរយៈពេលប៉ុន្មានថ្ងៃ?} & 12 & Trade Union \\ \midrule
22 & {\khmerfont តើហាមឃាត់ការរើសអើងដោយហេតុអ្វីក្នុងការចូលរួមជាសមាជិកសហជីព?} & 6 & Trade Union \\ \midrule
23 & {\khmerfont តើការចុះឈ្មោះអ្នកធានារ៉ាប់រងមានចែងក្នុងច្បាប់ដូចម្ដេច?} & 8 & Social Security \\ \midrule
24 & {\khmerfont តើការធានារ៉ាប់រងគ្រោះថ្នាក់ការងារ ផ្ដល់ប្រយោជន៍អ្វីខ្លះ?} & 37, 38 & Social Security \\ \midrule
25 & {\khmerfont តើរបបសន្តិសុខសង្គមមានបម្រើការប្រភេទណាខ្លះ?} & 5 & Social Security \\ \midrule
26 & {\khmerfont តើនរណាជាអ្នកកំណត់ប្រាក់ឈ្នួលអប្បបរមា?} & 8, 9 & Wages \& Compensation \\ \midrule
27 & {\khmerfont តើការរំលោភបំពានលើច្បាប់ប្រាក់ឈ្នួលអប្បបរមា ត្រូវទទួលទោសដូចម្ដេច?} & 22 & Wages \& Compensation \\ \midrule
28 & {\khmerfont តើការងារដោយបង្ខំ ត្រូវបានចែងដូចម្ដេចក្នុងច្បាប់ការងារ?} & 15, 16 & Employment Contract \\ \midrule
29 & {\khmerfont តើការព្យួរកិច្ចសន្យាការងារ មានហេតុអ្វីខ្លះ?} & 71 & Employment Contract \\ \midrule
30 & {\khmerfont តើការដកខ្លួនចេញពីសហជីព ប្រព្រឹត្តតាមដំណើរការដូចម្ដេច?} & 7 & Trade Union \\
\bottomrule
\caption{Complete Evaluation Dataset (30 Questions)}
\label{tab:appendix_dataset}
\end{longtable}

% ═══════════════════════════════════════════════════════════════════════
%  APPENDIX B — PER-QUESTION RETRIEVAL RESULTS
% ═══════════════════════════════════════════════════════════════════════
\section*{Appendix B: Per-Question Retrieval Results}
\addcontentsline{toc}{section}{Appendix B: Per-Question Retrieval Results}

\hs The following table reports, for each question in the final evaluation run, the
article numbers actually retrieved in the top-5, together with whether the expected
article was present (Retrieval Hit) and whether it was cited in the final answer
(Citation Hit).

\begin{longtable}{p{0.5cm} p{2cm} p{4.5cm} p{1.6cm} p{1.6cm}}
\toprule
\textbf{\#} & \textbf{Expected} & \textbf{Top-5 Retrieved} & \textbf{Retr. Hit} & \textbf{Cite Hit} \\
\midrule
\endhead
1  & 65       & 3, 5, 2, 8, 65        & Yes & Yes \\ \midrule
2  & 67       & 9, 10, 67, 1, 66      & Yes & Yes \\ \midrule
3  & 68       & 1, 10, 9, 3, 11       & No  & No  \\ \midrule
4  & 75       & 78, 1, 75, 2, 81      & Yes & Yes \\ \midrule
5  & 77       & 73, 1, 77, 2, 78      & Yes & Yes \\ \midrule
6  & 137      & 11, 137, 1, 141, 2    & Yes & Yes \\ \midrule
7  & 139      & 11, 139, 1, 144, 2    & Yes & Yes \\ \midrule
8  & 144      & 137, 139, 144, 2, 1   & Yes & Yes \\ \midrule
9  & 147, 146 & 147, 1, 146, 2, 148   & Yes & Yes \\ \midrule
10 & 166      & 11, 166, 1, 170, 2    & Yes & Yes \\ \midrule
11 & 166, 161 & 1, 11, 161, 170, 2    & Yes & Yes \\ \midrule
12 & 168      & 163, 1, 164, 2, 116   & No  & No  \\ \midrule
13 & 22       & 30, 20, 22, 19, 27    & Yes & Yes \\ \midrule
14 & 26       & 371, 1, 26, 2, 362    & Yes & Yes \\ \midrule
15 & 12       & 1, 12, 3, 15, 13      & Yes & Yes \\ \midrule
16 & 73       & 89, 1, 73, 2, 110     & Yes & Yes \\ \midrule
17 & 24       & 17, 20, 18, 19, 24    & Yes & Yes \\ \midrule
18 & 387, 389 & 98, 54, 300, 1, 385   & Yes & Yes \\ \midrule
19 & 10       & 49, 10, 56, 278, 64   & Yes & Yes \\ \midrule
20 & 5, 7     & 279, 1, 5, 2, 266     & Yes & Yes \\ \midrule
21 & 12       & 16, 1, 12, 2, 11      & Yes & Yes \\ \midrule
22 & 6        & 64, 1, 6, 2, 279      & Yes & Yes \\ \midrule
23 & 8        & 14, 82, 4, 7, 1       & No  & No  \\ \midrule
24 & 37, 38   & 66, 58, 1, 255, 2     & No  & No  \\ \midrule
25 & 5        & 7, 66, 13, 57, 1      & No  & No  \\ \midrule
26 & 8, 9     & 2, 11, 107, 1, 4      & No  & No  \\ \midrule
27 & 22       & 24, 1, 26, 2, 27      & Yes & Yes \\ \midrule
28 & 15, 16   & 1, 3, 2, 6, 15        & Yes & Yes \\ \midrule
29 & 71       & 71, 55, 72, 21, 95    & Yes & Yes \\ \midrule
30 & 7        & 7, 1, 293, 2, 43      & Yes & Yes \\
\bottomrule
\caption{Per-Question Retrieval Results (Final Run)}
\label{tab:appendix_results}
\end{longtable}

% ═══════════════════════════════════════════════════════════════════════
%  APPENDIX C — AGENT SYSTEM PROMPTS
% ═══════════════════════════════════════════════════════════════════════
\section*{Appendix C: Agent System Prompts}
\addcontentsline{toc}{section}{Appendix C: Agent System Prompts}

\hs This appendix reproduces the system prompts used by each language-model node in the
agent pipeline. They are presented exactly as used in the implementation.

\subsection*{C.1 Intent Classifier Prompt}
\begin{verbatim}
You are an intent classifier for a Cambodian legal chatbot.
Classify the user query into exactly one intent.

Intents:
- legal_qa        : general legal question requiring article retrieval
- article_lookup  : user asks for a specific article number
- definition      : user asks what a legal term means
- recent_news     : recent events, news, or current minimum-wage value
- greeting        : hello, thank you, chitchat, off-topic

Rules:
- If the query mentions a specific article number, use article_lookup
- Extract the article number as an Arabic numeral if present
- Questions about legal pay RATES, overtime, night work, or working
  hours are always legal_qa NOT recent_news
- Return ONLY valid JSON, no explanation

Output format:
{"intent": "<intent>", "article_number": "<number or null>"}
\end{verbatim}

\subsection*{C.2 Query Rewriter Prompt}
\begin{mdframed}
{\khmerfont\small
You are a Cambodian legal research assistant. Rewrite the user's question into an
optimised search query for retrieving relevant articles from Cambodian labour law
documents.\\[4pt]
The indexed documents cover exactly 4 laws:\\
1. ច្បាប់ស្តីពីការងារ (Labor Law) — employment contracts, wages, working hours, leave, termination\\
2. ច្បាប់ស្តីពីសហជីព (Trade Union Law) — union formation, registration, collective bargaining\\
3. ច្បាប់ស្តីពីរបបសន្តិសុខសង្គម (Social Security Law) — NSSF, work injury, health insurance, pension\\
4. ច្បាប់ស្តីពីប្រាក់ឈ្នួលអប្បបរមា (Minimum Wage Law) — minimum wage, wage council, violations\\[4pt]
Rules: expand informal terms into formal legal Khmer; include the law name only when
confident; do not add multiple law names unless the question spans them; keep it concise;
return only the rewritten query.
}
\end{mdframed}

\subsection*{C.3 Grounding Verification Prompt}
\begin{verbatim}
You are a relevance checker for a Cambodian legal RAG system.
Given a user question and retrieved legal article excerpts,
determine if the articles contain enough information to answer
the question.

Return ONLY valid JSON:
{"is_relevant": true} or {"is_relevant": false}
\end{verbatim}

\subsection*{C.4 Answer Generation Prompt}
\begin{mdframed}
{\khmerfont\small
អ្នកជាមេធាវីជំនាញច្បាប់កម្ពុជា។ ផ្តល់ចម្លើយជាផ្លូវការ ជាភាសាខ្មែរ។\\[4pt]
ក្បួនខ្នាត:\\
- ឆ្លើយសំណួរដោយផ្ទាល់ និងសង្ខេប — អតិបរមា ២០០ ពាក្យ\\
- ដកស្រង់មាត្រាច្បាប់ដែលពាក់ព័ន្ធដោយប្រើទម្រង់ [N] ដែល N គឺជាលេខយោងខាងក្រោម\\
- ប្រើភាសាច្បាប់ជំនាញ ប៉ុន្តែអាចយល់បាន\\
- ផ្តល់ចំណុចសំខាន់ ២-៣ ប៉ុណ្ណោះ; កុំរាយបញ្ជីវែង\\
- មិនបង្កើតព័ត៌មានដែលគ្មាននៅក្នុងឯកសារ\\[4pt]
ឆ្លើយតបជាអត្ថបទធម្មតា មិនមែន JSON។
}
\end{mdframed}

\subsection*{C.5 Citation Extraction Prompt}
\begin{verbatim}
Extract citations from this legal answer as JSON array only.
Return ONLY a valid JSON array, nothing else.

Format:
[
  {
    "ref_num": 1,
    "article_number": "95",
    "law_name": "...",
    "law_name_en": "Labor Law",
    "page_number": 30,
    "pdf_filename": "cambodian_labor_laws.pdf",
    "article_title": ""
  }
]

If no citations found, return: []
\end{verbatim}

% ═══════════════════════════════════════════════════════════════════════
%  APPENDIX D — KEY CODE LISTINGS
% ═══════════════════════════════════════════════════════════════════════
\section*{Appendix D: Key Code Listings}
\addcontentsline{toc}{section}{Appendix D: Key Code Listings}

\hs This appendix presents the core algorithms of the system. Listings are abbreviated
to the essential logic.

\subsection*{D.1 Reciprocal Rank Fusion (RRF)}
\begin{verbatim}
def _rrf(vec_hits, bm25_hits, k=60, top_n=5):
    """Fuse dense (FAISS) and sparse (BM25) ranked lists."""
    scores = {}
    cmap   = {}
    for rank, (_, chunk) in enumerate(vec_hits):
        cid = chunk["chunk_id"]
        scores[cid] = scores.get(cid, 0) + 1 / (k + rank + 1)
        cmap[cid]   = chunk
    for rank, chunk in enumerate(bm25_hits):
        cid = chunk["chunk_id"]
        scores[cid] = scores.get(cid, 0) + 1 / (k + rank + 1)
        cmap[cid]   = chunk
    ranked = sorted(scores, key=lambda x: scores[x], reverse=True)
    return [cmap[cid] for cid in ranked[:top_n]]
\end{verbatim}

\subsection*{D.2 Agent Pipeline Routing (LangGraph)}
\begin{verbatim}
def route_by_intent(state):
    intent = state["intent"]
    if intent == "article_lookup":
        return "direct_lookup"
    elif intent in ("legal_qa", "definition"):
        return "rewrite_query"
    else:                       # greeting / recent_news
        return "generate"

def route_after_grounding(state):
    if state["is_grounded"]:
        return "generate"
    if state["retry_count"] < 2:
        return "rewrite_query"   # retry with a refined query
    return "generate"            # give up retrying, generate anyway
\end{verbatim}

\subsection*{D.3 Two-Step Answer Generation with Citation Filter}
\begin{verbatim}
# Step 1 — generate the Khmer answer (no JSON pressure)
answer = generate_khmer_answer(query, context)

# Step 2 — extract citations in a separate call
citations = extract_citations(answer, chunks)

# Step 3 — keep only citations referenced inline as [N]
cited_nums = {int(n) for n in re.findall(r'\[(\d+)\]', answer)}
citations  = [c for c in citations if c["ref_num"] in cited_nums]
\end{verbatim}

% ═══════════════════════════════════════════════════════════════════════
%  APPENDIX E — DATABASE SCHEMA
% ═══════════════════════════════════════════════════════════════════════
\section*{Appendix E: Database Schema}
\addcontentsline{toc}{section}{Appendix E: Database Schema}

\hs The application uses a Neon serverless PostgreSQL database with the following schema.
The \texttt{chat\_sessions} and \texttt{chat\_messages} tables store conversation history,
while \texttt{anonymous\_usage} enforces the free-tier question limit for unauthenticated
visitors.

\begin{verbatim}
CREATE TABLE chat_sessions (
    session_id   TEXT PRIMARY KEY,
    user_id      TEXT,                 -- NULL = anonymous; Clerk user id otherwise
    created_at   TIMESTAMPTZ DEFAULT NOW(),
    updated_at   TIMESTAMPTZ DEFAULT NOW()
);

CREATE TABLE chat_messages (
    id           BIGSERIAL PRIMARY KEY,
    session_id   TEXT REFERENCES chat_sessions(session_id) ON DELETE CASCADE,
    role         TEXT NOT NULL CHECK (role IN ('user','assistant')),
    content      TEXT NOT NULL,
    citations    JSONB DEFAULT '[]',
    created_at   TIMESTAMPTZ DEFAULT NOW()
);

CREATE INDEX idx_messages_session ON chat_messages(session_id, created_at);
CREATE INDEX idx_sessions_user    ON chat_sessions(user_id, updated_at DESC);

CREATE TABLE anonymous_usage (
    anon_id        TEXT PRIMARY KEY,   -- httpOnly cookie identifier
    question_count INTEGER NOT NULL DEFAULT 0,
    created_at     TIMESTAMPTZ DEFAULT NOW(),
    updated_at     TIMESTAMPTZ DEFAULT NOW()
);
\end{verbatim}
